% CORRECTED VERSION
% Changes:
% - Fixed the algebraic error in Step 13 (the factor B was incorrectly cancelled).
% - Updated Theorem 1 to reflect the correct constant (includes factor B).
% - Added a brief remark after Lemma 5 clarifying the relation between block and full gradients.
% - Adjusted the Rate Derivation section to be consistent with the corrected bound.
% - Kept all original step annotations for easy reference.

\documentclass{article}
\usepackage{amsmath,amssymb,amsthm}
\usepackage{algorithm}
\usepackage{algorithmic}
\usepackage{enumitem}
\usepackage{hyperref}
\usepackage{color}
\newtheorem{theorem}{Theorem}
\newtheorem{lemma}{Lemma}
\newtheorem{assumption}{Assumption}
\newcommand{\E}{\mathbb{E}}
\newcommand{\R}{\mathbb{R}}
\newcommand{\norm}[1]{\left\|#1\right\|}
\newcommand{\grad}{\nabla}
\newcommand{\block}[1]{#1}
\newcommand{\iid}{\stackrel{\text{i.i.d.}}{\sim}}
\begin{document}

\section*{Algorithm Name}
\textbf{Randomized Block Coordinate Descent (RBCD)} for non‑convex smooth objectives.

\section*{Mathematical Formulation}
Let $f:\R^{d}\rightarrow\R$ be differentiable.  
Partition the coordinate vector $w\in\R^{d}$ into $B$ (possibly unequal) blocks
\[
w=\bigl(w^{(1)},w^{(2)},\dots ,w^{(B)}\bigr),\qquad 
w^{(i)}\in\R^{d_i},\;\; \sum_{i=1}^{B}d_i=d .
\]

At iteration $t\in\{0,1,\dots\}$

\begin{enumerate}[label=\arabic*)]
\item Sample a block index $i_t\in\{1,\dots ,B\}$ independently and uniformly:
\[
\Pr(i_t=i)=\frac1B,\qquad i=1,\dots ,B .
\]
\item Compute the partial (block) gradient
\[
g_t:=\grad_i f(w_t)\;\; \text{where } i=i_t .
\]
\item Update only the selected block with stepsize $\eta_t>0$:
\[
w_{t+1}^{(i)}=
\begin{cases}
w_{t}^{(i)}-\eta_t\,\grad_i f(w_t), & i=i_t,\\[4pt]
w_{t}^{(i)}, & i\neq i_t .
\end{cases}
\tag{1}
\]
\end{enumerate}

The full iterate $w_{t+1}$ is obtained by concatenating the updated block
and the untouched blocks.

\section*{Pseudocode}
\begin{algorithm}[H]
\caption{Randomized Block Coordinate Descent (RBCD)}
\begin{algorithmic}[1]
\REQUIRE Initial point $w_0\in\R^{d}$, stepsize schedule $\{\eta_t\}_{t\ge0}$, number of iterations $T$.
\FOR{$t=0$ \TO $T-1$}
    \STATE Sample block $i_t\sim\text{Uniform}\{1,\dots ,B\}$.
    \STATE Compute partial gradient $g_t=\grad_{i_t}f(w_t)$.
    \STATE Update the selected block:
    \[
    w_{t+1}^{(i_t)} = w_t^{(i_t)}-\eta_t g_t .
    \]
    \STATE Set $w_{t+1}^{(i)} = w_t^{(i)}$ for all $i\neq i_t$.
\ENDFOR
\RETURN $w_T$ (or optionally the averaged iterate $\bar w_T=\frac1T\sum_{t=0}^{T-1}w_t$).
\end{algorithmic}
\end{algorithm}

\section*{Dry Run Example}
Consider the one‑dimensional quadratic $f(w)=(w-3)^2$, which we treat as a single block ($B=1$).  
Take $w_0=0$, constant stepsize $\eta_t=\eta=0.1$, and run three iterations.

\begin{center}
\begin{tabular}{c|c|c|c}
$t$ & $w_t$ & $\grad f(w_t)=2(w_t-3)$ & $f(w_t)$\\\hline
0 & $0$ & $-6$ & $9$\\
1 & $w_1=0-0.1\cdot(-6)=0.6$ & $2(0.6-3)=-4.8$ & $(0.6-3)^2=5.76$\\
2 & $w_2=0.6-0.1\cdot(-4.8)=1.08$ & $2(1.08-3)=-3.84$ & $(1.08-3)^2=3.6864$\\
3 & $w_3=1.08-0.1\cdot(-3.84)=1.464$ & $2(1.464-3)=-3.072$ & $(1.464-3)^2=2.365$\\
\end{tabular}
\end{center}
The objective value decreases at each step, illustrating the descent
property of (1).

\newpage
\section*{Assumptions}
\begin{assumption}[Block‑wise $L$‑smoothness]\label{ass:smooth}
For each block $i\in\{1,\dots ,B\}$ there exists $L_i>0$ such that for any
$u,v\in\R^{d}$ differing only in block $i$,
\[
f(v)\le f(u)+\langle \grad_i f(u),\,v^{(i)}-u^{(i)}\rangle
+\frac{L_i}{2}\norm{v^{(i)}-u^{(i)}}^{2}.
\tag{2}
\]
\end{assumption}
Define $L_{\max}:=\max_{i}L_i$ and $L_{\mathrm{avg}}:=\frac1B\sum_{i=1}^{B}L_i$.

\begin{assumption}[Bounded below]\label{ass:lb}
The function $f$ is bounded from below, i.e.
\[
f_{\inf}:=\inf_{w\in\R^{d}}f(w)>-\infty .
\]
\end{assumption}

\begin{assumption}[Stepsize]\label{ass:stepsize}
The deterministic stepsize satisfies
\[
0<\eta_t\le \frac{1}{L_{\max}}\quad\text{for all }t\ge0 .
\tag{3}
\]
\end{assumption}

No convexity is required; $f$ may be non‑convex.

\section*{Main Convergence Theorem}
\begin{theorem}[Expected gradient norm decay]\label{thm:main}
Assume \ref{ass:smooth}--\ref{ass:stepsize} hold and let
$\eta_t\equiv\eta\le 1/L_{\max}$.  For any $T\ge1$,
\[
\frac1T\sum_{t=0}^{T-1}\E\!\bigl[\norm{\grad f(w_t)}^{2}\bigr]
\;\le\;
\frac{2\,B\,L_{\max}\bigl(f(w_0)-f_{\inf}\bigr)}{T}.
\tag{4}
\]
Consequently, there exists at least one iterate $t\le T-1$ with
\[
\E\!\bigl[\norm{\grad f(w_t)}^{2}\bigr]\le
\frac{2\,B\,L_{\max}\bigl(f(w_0)-f_{\inf}\bigr)}{T}.
\]
\end{theorem}
% NOTE: The factor $B$ is essential; it stems from the uniform
% random selection of a single block per iteration (see Step 13).

\section*{Theoretical Properties}
\begin{enumerate}[label=\arabic*)]
\item \textbf{Stability.} Because $\eta\le 1/L_{\max}$, the descent
inequality (2) guarantees that the expected function value is non‑increasing.
\item \textbf{Hyper‑parameter sensitivity.} The bound (4) scales linearly
with $B\,L_{\max}$; a larger Lipschitz constant forces a smaller admissible stepsize,
which slows the $1/T$ decay but preserves convergence.
\item \textbf{Comparison to full‑gradient SGD.} For the same stepsize,
RBCD reduces the per‑iteration computational cost from $O(d)$ to $O(d_i)$
while attaining the same $O(1/T)$ rate for the squared gradient norm up to the factor $B$ that reflects the block sampling.
\item \textbf{Special case: uniform Lipschitz constants.} If $L_i\equiv L$ for all $i$,
then $L_{\max}=L$ and (4) becomes $\frac{2BL(f(w_0)-f_{\inf})}{T}$.
\end{enumerate}

\section*{Preliminaries (Definitions and Lemmas)}
\begin{lemma}[Descent Lemma for a single block]\label{lem:descent}
Under Assumption~\ref{ass:smooth} and any stepsize $\eta\le 1/L_i$,
\[
f(w_{t+1})\le f(w_t)-\frac{\eta}{2}\norm{\grad_{i_t}f(w_t)}^{2}.
\tag{5}
\]
\end{lemma}
\begin{proof}
Apply (2) with $u=w_t$, $v=w_{t+1}$ and note that only block $i_t$ changes:
\[
f(w_{t+1})\le f(w_t)-\eta\langle \grad_{i_t}f(w_t),\grad_{i_t}f(w_t)\rangle
+\frac{L_{i_t}}{2}\eta^{2}\norm{\grad_{i_t}f(w_t)}^{2}.
\]
Since $\eta\le 1/L_{i_t}$, we have $L_{i_t}\eta\le1$, thus
\[
-\eta+\frac{L_{i_t}}{2}\eta^{2}
=-\frac{\eta}{2}\bigl(2-L_{i_t}\eta\bigr)\le -\frac{\eta}{2}.
\]
Rearranging yields (5). \hfill $\square$
\end{proof}
% ADDED: Clarify the relation between block and full gradients.
\begin{remark}
For any $w$, the full gradient decomposes as
\[
\grad f(w)=\bigl(\grad_1 f(w),\dots,\grad_B f(w)\bigr),\qquad
\norm{\grad f(w)}^{2}=\sum_{i=1}^{B}\norm{\grad_i f(w)}^{2}.
\tag{5a}
\]
\end{remark}

\section*{Main Proof (Step‑by‑Step Derivation)}
We prove Theorem~\ref{thm:main}.  All expectations are taken w.r.t. the
random choices $\{i_0,\dots ,i_{t}\}$ conditioned on the past iterates.

\begin{enumerate}[label=[Step \arabic*]]
\item \textbf{One‑step expected descent.}  
Take expectation of Lemma~\ref{lem:descent} conditioned on $w_t$:
\[
\E\!\bigl[f(w_{t+1})\mid w_t\bigr]
\le f(w_t)-\frac{\eta}{2}\E\!\bigl[\norm{\grad_{i_t}f(w_t)}^{2}\mid w_t\bigr].
\tag{6}
\]

\item \textbf{Relate block gradient to full gradient.}  
Because $i_t$ is uniform,
\[
\E\!\bigl[\norm{\grad_{i_t}f(w_t)}^{2}\mid w_t\bigr]
=\frac1B\sum_{i=1}^{B}\norm{\grad_i f(w_t)}^{2}
=\frac1B\norm{\grad f(w_t)}^{2}.
\tag{7}
\]

\item \textbf{Combine (6) and (7).}
\[
\E\!\bigl[f(w_{t+1})\mid w_t\bigr]
\le f(w_t)-\frac{\eta}{2B}\norm{\grad f(w_t)}^{2}.
\tag{8}
\]

\item \textbf{Take full expectation.}  
Denote $\Phi_t:=\E[f(w_t)]$.  Using the law of total expectation on (8):
\[
\Phi_{t+1}\le \Phi_t-\frac{\eta}{2B}\E\!\bigl[\norm{\grad f(w_t)}^{2}\bigr].
\tag{9}
\]

\item \textbf{Telescoping sum.}  
Sum (9) from $t=0$ to $T-1$:
\[
\Phi_T\le \Phi_0-\frac{\eta}{2B}\sum_{t=0}^{T-1}
\E\!\bigl[\norm{\grad f(w_t)}^{2}\bigr].
\tag{10}
\]

\item \textbf{Lower bound $\Phi_T$.}  
By Assumption~\ref{ass:lb}, $\Phi_T\ge f_{\inf}$.  Rearranging (10):
\[
\sum_{t=0}^{T-1}\E\!\bigl[\norm{\grad f(w_t)}^{2}\bigr]
\le \frac{2B}{\eta}\bigl(\Phi_0-f_{\inf}\bigr).
\tag{11}
\]

\item \textbf{Insert stepsize bound.}  
Choose the maximal admissible constant stepsize $\eta=1/L_{\max}$ (allowed by
Assumption~\ref{ass:stepsize}).  Then (11) becomes
\[
\sum_{t=0}^{T-1}\E\!\bigl[\norm{\grad f(w_t)}^{2}\bigr]
\le 2B L_{\max}\bigl(f(w_0)-f_{\inf}\bigr).
\tag{12}
\]

\item \textbf{Average over $T$ iterations.}
Divide (12) by $T$:
\[
\frac1T\sum_{t=0}^{T-1}\E\!\bigl[\norm{\grad f(w_t)}^{2}\bigr]
\le \frac{2\,B\,L_{\max}\bigl(f(w_0)-f_{\inf}\bigr)}{T}.
\tag{13}
\]
% CORRECTED: The factor $B$ is retained; it cannot be cancelled.
Equation (13) is exactly the bound (4) stated in Theorem~\ref{thm:main}.
\hfill $\square$
\end{enumerate}

\section*{Rate Derivation (With Constants)}
From (13) we read the explicit convergence constant:
\[
C:=2\,B\,L_{\max}\bigl(f(w_0)-f_{\inf}\bigr).
\]
Thus the expected squared gradient norm decays as $C/T$.
If a diminishing stepsize $\eta_t=\frac{c}{t+1}$ with $c\le 1/L_{\max}$ is used,
the same derivation yields an $O(\log T/T)$ bound; however the constant
stepsize provides the clean $O(1/T)$ rate.

\section*{Special Cases}
\begin{itemize}
\item \textbf{Uniform Lipschitz constants.} If $L_i\equiv L$, then $L_{\max}=L$ and (4) reduces to
\[
\frac1T\sum_{t=0}^{T-1}\E\!\bigl[\norm{\grad f(w_t)}^{2}\bigr]\le
\frac{2\,B\,L\bigl(f(w_0)-f_{\inf}\bigr)}{T}.
\]
\item \textbf{Single block ($B=1$).} The algorithm coincides with standard gradient descent,
and (4) recovers the classical $O(1/T)$ rate for smooth non‑convex optimization.
\item \textbf{Large‑scale sparse problems.} When each block contains only a few coordinates,
the per‑iteration cost is dramatically reduced while the convergence guarantee remains unchanged up to the factor $B$.
\end{itemize}

\section*{Discussion}
The proof hinges only on block‑wise smoothness and a stepsize not exceeding the
largest block Lipschitz constant.  No convexity or variance assumptions are needed,
which makes the result applicable to a broad class of non‑convex machine‑learning
problems (e.g., training of deep networks with coordinate‑wise updates, matrix factorization,
and structured sparsity models).  The $O(1/T)$ decay matches the optimal rate for
first‑order methods on smooth non‑convex functions; the extra factor $B$ simply
reflects that only one out of $B$ blocks is updated per iteration.

\end{document}
% ===== CORRECTION SUMMARY =====
% 1. [Step 13] Issue: Incorrect cancellation of the factor B. Fixed by retaining B in the final bound.
% 2. Theorem 1 was updated to include the factor B, making it consistent with the corrected derivation.
% 3. Added a remark (5a) clarifying the relationship between block and full gradients.
% 4. Adjusted the Rate Derivation section to reflect the corrected constant.